The ASR and RAG experiments expose a common systems
trade-off: increasing computational work or resource
use does not necessarily improve the desired outcome.
For ASR, the custom CUDA implementation passed
correctness tests but was substantially slower than
the PyTorch baseline. Controlled substitutions and
profiling identified GEMM as the dominant bottleneck.
The tile-size experiments further showed that a
single kernel configuration is poorly suited to
both large encoder matrices and the one-row
\texttt{lm\_head} workload.

For RAG, more retrieval recall need not produce
better answers: recall records evidence coverage,
not whether the context is concise or coherent.
Additional documents can introduce irrelevant
evidence or dilute useful material, while the
generator may still misread, omit, or fail to
combine relevant facts. The oracle result is
consistent with evidence availability mattering,
but its context is not matched in document count
or token length and does not isolate retrieval.

These conclusions are bounded by the experimental
settings. ASR timings used a single input on an
H200 MIG partition, with limited repeated runs;
the hybrid's small measured advantage over
PyTorch requires replication. RAG used 125
development queries and an automated judge;
its latency subset covered 20 queries. The
RTX 2080 Ti experiments conducted separately
are not directly comparable to H200 MIG timings.

Future work would use shape-specialised GEMM
kernels, improved attention tiling and broader
ASR workload testing. For RAG, stronger
retrieval and evidence selection, repeated
generations, independent human assessment
and evaluation on a held-out query set would
provide stronger evidence of quality--cost
trade-offs. Overall, the results favour
workload-specific optimisation and controlled
end-to-end measurement over isolated kernel
speed or context-size increases.